\makeatletter
\@ifundefined{keepfig}{\newenvironment{keepfig}{\par\edef\keepfig@pi{\the\parindent}\noindent
  \begin{minipage}{\linewidth}\parindent=\keepfig@pi\relax}{\end{minipage}\par}}{}
\makeatother
\section*{Упражнения}
\addcontentsline{toc}{section}{Упражнения}

\begin{exercise}[1-1]
Пусть $E=\Q[\alpha]$, где $\alpha^3-\alpha^2+\alpha+2=0$. Запишите $(\alpha^2+\alpha+1)(\alpha^2-\alpha)$ и $(\alpha-1)^{-1}$ в виде
$a\alpha^2+b\alpha+c$, $a,b,c\in\Q$.
\end{exercise}

\noindent\emph{Решение.} Многочлен $f=X^3-X^2+X+2$ неприводим над $\Q$: многочлен степени $3$ приводим лишь при
наличии линейного множителя, т.\,е.\ рационального корня, а рациональные корни унитарного целочисленного многочлена~---
целые делители свободного члена (\rf{1.11}); но $f(1)=3$, $f(-1)=-1$, $f(2)=8$, $f(-2)=-12$. Поэтому
$E\cong\Q[X]/(f)$~--- поле степени $3$ с базисом $1,\alpha,\alpha^2$ (\rf{1.25}), и вычисления сводятся к понижению степени с
помощью $\alpha^3=\alpha^2-\alpha-2$. Отсюда $\alpha^4=\alpha^3-\alpha^2-2\alpha=-3\alpha-2$ и
\[
(\alpha^2+\alpha+1)(\alpha^2-\alpha)=\alpha^4-\alpha=-4\alpha-2 .
\]
Для обратного элемента разделим $f$ на $X-1$ с остатком: $f=(X-1)(X^2+1)+3$. Подставив $\alpha$, получим
$(\alpha-1)(\alpha^2+1)=-3$, т.\,е.
\[
(\alpha-1)^{-1}=-\tfrac13\alpha^2-\tfrac13 .
\]
(В общем случае обратный элемент находится алгоритмом Евклида, как в построении поля $F[X]/(f)$ и в примере \rf{1.28};
здесь хватило одного деления.)

\begin{exercise}[1-2]
Найдите $[\Q(\sqrt2,\sqrt3):\Q]$.
\end{exercise}

\noindent\emph{Решение.} Ответ: $4$. Имеем $\Q\subset\Q(\sqrt2)\subset\Q(\sqrt2,\sqrt3)=\Q(\sqrt2)[\sqrt3]$. Первая ступень имеет
степень $2$ ($X^2-2$ неприводим: $\sqrt2\notin\Q$). Вторая имеет степень $2$, если $\sqrt3\notin\Q(\sqrt2)$. Пусть
$\sqrt3=a+b\sqrt2$, $a,b\in\Q$. Тогда $3=a^2+2b^2+2ab\sqrt2$, и так как $1,\sqrt2$ линейно независимы над $\Q$, то
$ab=0$ и $a^2+2b^2=3$. При $b=0$ получаем $a^2=3$, а при $a=0$~--- $(2b)^2=6$; ни $3$, ни $6$ не квадраты рациональных
чисел. Значит, $\sqrt3\notin\Q(\sqrt2)$ (подробнее и в общем виде~--- \rf{1.32a}), и по \rf{1.20}
$[\Q(\sqrt2,\sqrt3):\Q]=2\cdot2=4$ с базисом $1,\sqrt2,\sqrt3,\sqrt6$. Заметим, что $\Q(\sqrt2,\sqrt3)=\Q(\sqrt2+\sqrt3)$:
так как $(\sqrt3+\sqrt2)(\sqrt3-\sqrt2)=1$, поле $\Q(\sqrt2+\sqrt3)$ содержит $\sqrt3-\sqrt2$, а значит, и
$\sqrt3=\frac12\bigl((\sqrt3+\sqrt2)+(\sqrt3-\sqrt2)\bigr)$, $\sqrt2=\frac12\bigl((\sqrt3+\sqrt2)-(\sqrt3-\sqrt2)\bigr)$. Число
$\sqrt2+\sqrt3$~--- корень многочлена $X^4-10X^2+1=(X^2-5)^2-24$ (так как $(\sqrt2+\sqrt3)^2-5=2\sqrt6$), а этот
многочлен неприводим над $\Q$, но приводим по модулю каждого простого числа (\rf{1.18c}); значит, это минимальный
многочлен $\sqrt2+\sqrt3$ (\rf{1.27c}).

\begin{exercise}[1-3]
Пусть $F$~--- поле, $f(X)\in F[X]$.
\begin{enumerate}[label=(\alph*)]
\item Покажите, что для каждого $a\in F$ существует многочлен $q(X)\in F[X]$ такой, что $f(X)=q(X)(X-a)+f(a)$.
\item Выведите, что $f(a)=0$ тогда и только тогда, когда $(X-a)\mid f(X)$.
\item Выведите, что ненулевой многочлен $f(X)$ имеет не более $\deg f$ корней.
\item Пусть $G$~--- конечная абелева группа. Покажите: если для каждого делителя $m$ порядка $(G:1)$ в $G$ не более $m$
элементов, порядок которых делит $m$, то $G$ циклическая.
\item Выведите, что каждая конечная подгруппа группы $F^\times$ циклическая.
\end{enumerate}
\end{exercise}

\noindent\emph{Решение.} (a) Если $f=\sum c_kX^k$, то $f(X)-f(a)=\sum c_k(X^k-a^k)$, и каждое слагаемое делится на $X-a$:
$X^k-a^k=(X-a)(X^{k-1}+aX^{k-2}+\dots+a^{k-1})$. (Иначе: разделим $f$ на $X-a$ с остатком; остаток~--- константа $r$, и,
подставив $X=a$, получаем $r=f(a)$; это теорема Безу.) Ещё один способ~--- сдвиг переменной. Пусть $\deg f\le n$;
положим $g(X)=f(X+a)=\sum_{k=0}^ng_kX^k$. Тогда $g_0=g(0)=f(a)$ и
\[
f(X)=g(X-a)=f(a)+(X-a)\bigl(g_1+g_2(X-a)+\dots+g_n(X-a)^{n-1}\bigr)
\]
(коэффициенты $g_k$~--- это $(D^{(k)}f)(a)$ из формулы Тейлора \rf{1.7a}).

(b) Следует из (a): $f(X)=q(X)(X-a)+f(a)$, и $X-a$ делит $f$ тогда и только тогда, когда делит константу $f(a)$, т.\,е.\
когда $f(a)=0$.

(c) Это доказано в \rf{1.7} индукцией по $\deg f$ с помощью (b): если $f=(X-a)q$ и $b\ne a$~--- ещё один корень, то
$0=(b-a)q(b)$, откуда $q(b)=0$, так как в поле нет делителей нуля. Там же пример с $\Z/8\Z$, показывающий, что
отсутствие делителей нуля существенно.

(d) Уточним, что утверждается. Пусть $n=(G:1)$. Условие: для каждого $m\mid n$ уравнение $x^m=1$ имеет в $G$ не
более $m$ решений (элемент удовлетворяет $x^m=1$ тогда и только тогда, когда его порядок делит $m$). Вывод: в $G$ есть
элемент порядка $n$. Для циклической группы $G=\langle g\rangle$ порядка $n$ условие выполнено с равенством: при $m\mid n$
решения $x^m=1$~--- это ровно $m$ элементов подгруппы $\langle g^{n/m}\rangle$; например, в $\Z/6\Z$ при $m=1,2,3,6$
уравнению $mx=0$ удовлетворяют соответственно $\{0\}$, $\{0,3\}$, $\{0,2,4\}$ и все шесть элементов. Так что условие
характеризует циклические группы. Нециклическая группа его нарушает: например, в $\Z/4\Z\times\Z/2\Z$ ($n=8$) при $m=2$ уравнению $2x=0$
удовлетворяют четыре элемента $(0,0),(2,0),(0,1),(2,1)$.

\emph{Доказательство через экспоненту.} Пусть $e$~--- \emph{экспонента} $G$, т.\,е.\ наименьшее общее кратное порядков
всех элементов; $e\mid n$ по теореме Лагранжа. Ключевой факт: \emph{в конечной абелевой группе есть элемент порядка $e$};
иными словами, порядок каждого элемента делит наибольший из порядков. Докажем его. Если $\operatorname{ord} a=r$, $\operatorname{ord} b=s$,
$\gcd(r,s)=1$ и $ab=ba$, то $\operatorname{ord}(ab)=rs$: ясно, что $(ab)^{rs}=1$, а если $(ab)^t=1$, то $a^t=b^{-t}$ лежит в
$\langle a\rangle\cap\langle b\rangle$, порядок которой делит $r$ и $s$, т.\,е.\ равен $1$; поэтому $a^t=b^t=1$, $r\mid t$,
$s\mid t$, $rs\mid t$. Пусть $e=p_1^{k_1}\cdots p_l^{k_l}$. Для каждого $i$ найдётся $x_i\in G$, порядок которого делится
на $p_i^{k_i}$, и тогда $y_i=x_i^{\operatorname{ord}(x_i)/p_i^{k_i}}$ имеет порядок ровно $p_i^{k_i}$. По доказанному элемент
$y_1\cdots y_l$ имеет порядок $p_1^{k_1}\cdots p_l^{k_l}=e$.

Теперь все $n$ элементов $G$ удовлетворяют уравнению $x^e=1$, а по условию (с $m=e$) таких элементов не больше $e$.
Значит, $n\le e$; вместе с $e\mid n$ это даёт $e=n$, и элемент порядка $e$ порождает $G$.

Для неабелевых групп ключевой факт неверен: в $S_3$ наибольший порядок элемента равен $3$, но транспозиции имеют
порядок $2$. Тем не менее само утверждение (d) верно и без абелевости. \emph{Доказательство подсчётом.} Обозначим через
$\psi(d)$ число элементов порядка $d$ в $G$, через $\varphi$~--- функцию Эйлера. Если $\psi(d)>0$ и $\operatorname{ord} x=d$, то все $d$
элементов группы $\langle x\rangle$ удовлетворяют $y^d=1$, а по условию таких элементов в $G$ не больше $d$; значит, все
элементы порядка $d$ лежат в $\langle x\rangle$ и являются её образующими, и $\psi(d)=\varphi(d)$. Итак, $\psi(d)\le\varphi(d)$
для всех $d\mid n$, и
\[
n=\sum_{d\mid n}\psi(d)\le\sum_{d\mid n}\varphi(d)=n
\]
(последнее равенство~--- подсчёт элементов циклической группы $\Z/n\Z$ по их порядкам: элементов порядка $d$ в ней
ровно $\varphi(d)$). Значит, всюду равенства, и $\psi(n)=\varphi(n)\ge1$.

Для абелевых групп можно рассуждать и через теорему о строении конечных абелевых групп: если $G$ не циклическая, то в
разложении $G$ в произведение циклических групп примарных порядков встречаются два множителя $\Z/p^\alpha\Z$ и
$\Z/p^\beta\Z$ с одним и тем же простым $p$ (иначе по китайской теореме об остатках $G$ циклическая), и подгруппа
$\Z/p\Z\times\Z/p\Z$ даёт $p^2>p$ решений уравнения $x^p=1$. Используется только существование такого
разложения; его единственность не нужна. Вариант того же рассуждения~--- через экспоненту: если $G\cong\bigoplus_{i=1}^r\Z/p_i^{k_i}\Z$
(простые $p_i$ могут повторяться), то $n=\prod_ip_i^{k_i}$, а экспонента $e=\operatorname{\textup{НОК}}_i(p_i^{k_i})$; по
условию (с $m=e$) $n\le e$, т.\,е.\ $\operatorname{\textup{НОК}}_i(p_i^{k_i})=\prod_ip_i^{k_i}$, а это возможно лишь при попарно
различных $p_i$, и тогда $G\cong\Z/n\Z$ по китайской теореме об остатках.

(e) Пусть $G\subset F^\times$~--- конечная подгруппа. Решения уравнения $x^m=1$ в $G$~--- корни многочлена $X^m-1$, а их
в поле $F$ по (c) не больше $m$. По (d) группа $G$ циклическая. В частности, циклична мультипликативная группа
конечного поля; например, $\F_7^\times=\langle3\rangle$: степени $3$ по модулю $7$ равны $3,2,6,4,5,1$. Отсюда также:
корни степени $n$ из единицы в любом поле образуют циклическую группу, а конечное поле $E$ порождается над любым своим
подполем одним элементом~--- образующей группы $E^\times$.

\begin{exercise}[1-4]
Покажите, что с помощью циркуля, линейки и трисектора можно построить правильный $7$-угольник.
\end{exercise}

\noindent\emph{Решение.} Как в рассуждении после \rf{1.42}, достаточно построить $\alpha=\cos\frac{2\pi}7$: перпендикуляр к оси $x$ в точке
$(\alpha,0)$ пересекает единичную окружность в вершине $P_1=(\cos\frac{2\pi}7,\sin\frac{2\pi}7)$, и сторона $P_0P_1$,
$P_0=(1,0)$, откладывается циркулем по окружности.

Какие углы можно строить? Прямой угол строится (перпендикуляр), и любой построенный угол можно разделить пополам
(биссектриса). Алгебраически это видно из формулы $\cos2x=2\cos^2x-1$: при $0\le\theta\le\pi$
\[
\cos\frac\theta2=\sqrt{\frac{1+\cos\theta}2},
\]
т.\,е.\ $\cos\frac\theta2$ получается из $\cos\theta$ арифметическими действиями и одним квадратным корнем (\rf{1.36}).
Поэтому строятся все углы $\frac{2\pi}{2^n}$, а с трисектором~--- ещё и углы $\frac{2\pi}{3^m}$, а также их суммы и
разности (\rf{1.35a}(d)). Но угол $\frac{2\pi}7$ так получить нельзя: делением на два и на три, сложением и вычитанием
получаются лишь углы $2\pi\frac a{2^n3^m}$, $a\in\Z$, а $\frac17$ не имеет такого вида ($7\nmid2^n3^m$). Нужна алгебра:
по \rf{1.42a} с трисектором строятся корни многочленов $4Y^3-3Y-\beta$, $|\beta|\le1$, и надо свести к такому уравнению
уравнение для $\alpha$.

\emph{Шаг 1: уравнение для $\alpha$.} Пусть $\zeta=e^{2\pi i/7}$ и $t=\zeta+\zeta^{-1}=2\alpha$. Разделив
$1+\zeta+\dots+\zeta^6=0$ на $\zeta^3$, получим $(\zeta^3+\zeta^{-3})+(\zeta^2+\zeta^{-2})+(\zeta+\zeta^{-1})+1=0$, а
$\zeta^2+\zeta^{-2}=t^2-2$, $\zeta^3+\zeta^{-3}=t^3-3t$. Значит,
\[
t^3+t^2-2t-1=0,\qquad\text{т.\,е.}\qquad 8\alpha^3+4\alpha^2-4\alpha-1=0 .
\]
Многочлен $X^3+X^2-2X-1$ не имеет рациональных корней (в точках $\pm1$ он равен $-1$ и $1$) и потому неприводим;
$[\Q[\alpha]:\Q]=3$, и без трисектора $\alpha$ не построить (\rf{1.38}). Тот же вывод для $\zeta^2,\zeta^3$ вместо $\zeta$
(это тоже корни многочлена $1+X+\dots+X^6$: они отличны от $1$, так как $7\nmid2$, $7\nmid3$, а их седьмые степени
равны $1$) показывает, что корни многочлена $\mu(X)=X^3+\frac12X^2-\frac12X-\frac18$~--- это $\cos\frac{2\pi}7$, $\cos\frac{4\pi}7$,
$\cos\frac{6\pi}7$; все они вещественны (ср.\ \rf{3.22}, где разобрано поле $\Q[\zeta]$ и тот же многочлен
$X^3+X^2-2X-1$). Вещественность корней видна и без тригонометрии: многочлен $8\mu=8X^3+4X^2-4X-1$ принимает в
точках $-1$, $-\frac12$, $0$, $1$ значения $-1$, $1$, $-1$, $7$, поэтому у него есть корень в каждом из интервалов
$(-1,-\frac12)$, $(-\frac12,0)$, $(0,1)$; это $\cos\frac{6\pi}7<\cos\frac{4\pi}7<\cos\frac{2\pi}7$.

\begin{center}
\begin{keepfig}
\centering
\begin{tikzpicture}[every node/.style={font=\small}]
\begin{scope}[scale=1.5]
\draw[->,gray] (-1.25,0) -- (1.3,0);
\draw[gray] (0,0) circle (1);
\foreach \k in {1,...,3}{
  \draw[cnew!60,dashed] ({cos(360*\k/7)},{sin(360*\k/7)}) -- ({cos(360*\k/7)},{-sin(360*\k/7)});
  \fill[cnew] ({cos(360*\k/7)},0) circle (1pt);}
\fill[cold] (1,0) circle (1.1pt) node[below right=-1pt] {$1$};
\foreach \k/\lab/\pos in {1/\zeta/above right,2/\zeta^2/above left,3/\zeta^3/left,4/\zeta^4/left,5/\zeta^5/below left,6/\zeta^6/below right}{
  \fill ({cos(360*\k/7)},{sin(360*\k/7)}) circle (1.1pt) node[\pos=-1pt] {$\lab$};}
\end{scope}
\begin{scope}[xshift=5.9cm,yscale=0.42,xscale=2.2]
\draw[->,gray] (-1.2,0) -- (1.2,0) node[right] {$x$};
\draw[->,gray] (0,-2) -- (0,7.4);
\draw[cold,thick,domain=-1.08:1.02,samples=80] plot (\x,{8*\x*\x*\x+4*\x*\x-4*\x-1});
\foreach \x/\y/\p in {-1/-1/left,-0.5/1/above,0/-1/below left,1/7/right}{\node[circle,fill=cold,inner sep=1pt,label={[cold,font=\footnotesize]\p:$\y$}] at (\x,\y) {};}
\foreach \k in {1,2,3}{\node[circle,fill=cnew,inner sep=1.1pt] at ({cos(360*\k/7)},0) {};}
\draw (-1,0.15) -- (-1,-0.15); \draw (-0.5,0.15) -- (-0.5,-0.15) node[below left=-2pt] {$-\frac12$}; \draw (1,0.15) -- (1,-0.15) node[below] {$1$};
\end{scope}
\end{tikzpicture}
\figcap{Слева: корни $\zeta^k$ многочлена $X^7-1$; точки $\zeta^k$ и $\zeta^{-k}$ симметричны относительно оси $x$, и их
общая проекция~--- $\cos\frac{2\pi k}7=\frac{\zeta^k+\zeta^{-k}}2$. Справа: график $8x^3+4x^2-4x-1$; смена знака на
отрезках $[-1,-\frac12]$, $[-\frac12,0]$, $[0,1]$ даёт три вещественных корня $\cos\frac{6\pi}7<\cos\frac{4\pi}7<\cos\frac{2\pi}7$.}
\end{keepfig}
\end{center}

\emph{Шаг 2: тригонометрическое решение кубического уравнения.} Пусть $F\subset\R$~--- поле и $g=X^3+aX^2+bX+c\in F[X]$
имеет три вещественных корня (с учётом кратности). Подстановка $X=T-\frac a3$ убирает квадратичный член:
$g(T-\frac a3)=T^3+pT+q$, где $p=b-\frac{a^2}3$ и $q=\frac{2a^3}{27}-\frac{ab}3+c$ лежат в $F$. Для корней $t_1,t_2,t_3$ имеем $\sum t_i=0$, $\sum_{i<j}t_it_j=p$, поэтому
$\sum t_i^2=-2p$; значит, $p\le0$, причём $p=0$ (т.\,е.\ $a^2=3b$) лишь при $t_1=t_2=t_3=0$ (тогда все корни $g$ равны
$-\frac a3\in F$).
Пусть $p<0$. Положим $d=-\frac p3>0$, $k=2\sqrt d$ (так что $k^2=-\frac{4p}3$) и $T=kY$. Умножив
$k^3Y^3+pkY+q=0$ на $\frac4{k^3}$ и учтя $\frac{4p}{k^2}=-3$, получим
\[
4Y^3-3Y=\beta,\qquad \beta=-\frac{4q}{k^3}\in F[\sqrt d].
\]
При этом $|\beta|\le1$: так как $\beta^2=\frac{16q^2}{k^6}=-\frac{27q^2}{4p^3}$ и $p<0$, неравенство $\beta^2\le1$
равносильно $\Delta:=-4p^3-27q^2\ge0$, а $\Delta$~--- дискриминант:
\[
\Delta=\prod_{i<j}(t_i-t_j)^2\ge0,
\]
поскольку корни вещественны. Формулу для $\Delta$ проще всего получить через определитель Вандермонда $V=(t_j^{\,i})$,
$0\le i\le2$, $1\le j\le3$: $\prod_{i<j}(t_i-t_j)^2=(\det V)^2=\det(VV^{\mathsf T})$, а $VV^{\mathsf T}$~--- матрица степенных
сумм $s_k=\sum t_i^k$. Из $t_i^3=-pt_i-q$ получаем $s_0=3$, $s_1=0$, $s_2=-2p$, $s_3=-ps_1-3q=-3q$, $s_4=-ps_2-qs_1=2p^2$, и
\[
\Delta=\det\begin{pmatrix}3&0&-2p\\0&-2p&-3q\\-2p&-3q&2p^2\end{pmatrix}=3(-4p^3-9q^2)+(-2p)\cdot(-4p^2)=-4p^3-27q^2 .
\]
Итак, корни $g$ равны $-\frac a3+ky$, где $y$ пробегает корни $4Y^3-3Y-\beta$, т.\,е.\ числа
$\cos\frac{\varphi+2\pi j}3$, $j=0,1,2$, где $\varphi=\arccos\beta$. Они строятся циркулем, линейкой и трисектором, если
построимы коэффициенты $g$. (Это <<неприводимый случай>> формулы Кардано: при трёх вещественных корнях она требует
кубических корней из невещественных чисел, а тригонометрическая подстановка заменяет их трисекцией угла.)

\emph{Шаг 3: вычисление.} Для $\mu$ имеем $a=\frac12$; подстановка $X=T-\frac16$ даёт
\[
\mu\bigl(T-\tfrac16\bigr)=T^3-\tfrac7{12}T-\tfrac7{216},\qquad p=-\tfrac7{12},\quad q=-\tfrac7{216}.
\]
Далее $d=\frac7{36}$, $k=\frac{\sqrt7}3$, $k^3=\frac{7\sqrt7}{27}$, $\beta=\frac{28}{216}\cdot\frac{27}{7\sqrt7}=\frac1{2\sqrt7}=\frac{\sqrt7}{14}$,
и $\Delta=\frac{1372-49}{1728}=\frac{49}{64}>0$. Пусть $\varphi=\arccos\frac1{2\sqrt7}\in(0,\frac\pi2)$. Тогда
$y_0=\cos\frac\varphi3>0$, а $\cos\frac{\varphi+2\pi}3$ и $\cos\frac{\varphi+4\pi}3$ отрицательны (углы лежат в
$(\frac{2\pi}3,\frac{5\pi}6)$ и $(\frac{4\pi}3,\frac{3\pi}2)$). Отображение $y\mapsto-\frac16+\frac{\sqrt7}3y$ возрастает, а
наибольший корень $\mu$~--- это $\cos\frac{2\pi}7$ (два других отрицательны). Поэтому
\[
\cos\frac{2\pi}7=-\frac16+\frac{\sqrt7}3\cos\Bigl(\frac13\arccos\frac1{2\sqrt7}\Bigr)
\]
(численно обе части равны $0{,}6234898\ldots$), и два других корня $\mu$ отвечают $j=1,2$.

\emph{Шаг 4: построение.}
\begin{enumerate}
\item Построим единичную окружность с центром $O$ и угол $\varphi$ с $\cos\varphi=\frac1{2\sqrt7}$: это угол прямоугольного
треугольника с катетами $1$ и $3\sqrt3$, прилежащий к катету $1$ (гипотенуза $\sqrt{1+27}=2\sqrt7$, $\operatorname{tg}\varphi=3\sqrt3$;
$\varphi\approx79{,}11^\circ$). Числа $\sqrt7$ и $3\sqrt3=\sqrt{27}$ построимы по \rf{1.36}.
\item Трисектором разделим $\varphi$ на три части: $\psi=\frac\varphi3\approx26{,}37^\circ$.
\item На стороне угла $\psi$ отложим $OM=\frac{\sqrt7}3\approx0{,}882$ и опустим перпендикуляр $MH$ на ось $x$:
$OH=\frac{\sqrt7}3\cos\psi\approx0{,}790$.
\item Отложим от $H$ к $O$ отрезок $\frac16$: получим точку $(\alpha,0)$, $\alpha\approx0{,}623$.
\item Перпендикуляр к оси $x$ в точке $(\alpha,0)$ пересечёт единичную окружность в вершине $P_1$; откладывая хорду
$P_0P_1$ по окружности, получим остальные вершины.
\end{enumerate}

\begin{center}
\begin{keepfig}
\centering
\begin{tikzpicture}[scale=0.62,baseline={(0,2.6)},every node/.style={font=\small}]
\coordinate (B) at (0,0); \coordinate (A) at (1,0); \coordinate (C) at (1,5.196);
\draw[cold,thick] (B) -- node[below] {$1$} (A) -- node[right] {$3\sqrt3$} (C);
\draw[cnew,thick] (B) -- node[above left=-1pt] {$2\sqrt7$} (C);
\draw (A) ++(-0.3,0) -- ++(0,0.3) -- ++(0.3,0);
\draw[cnew] (0.6,0) arc (0:79.1:0.6); \node[cnew] at (0.55,0.95) {$\varphi$};
\fill (B) circle (2pt) node[below left] {$B$}; \fill (A) circle (2pt) node[below right] {$A$};
\fill (C) circle (2pt) node[above] {$C$};
\end{tikzpicture}\qquad\qquad
\begin{tikzpicture}[scale=3.3,baseline={(0,0)},every node/.style={font=\small}]
\pgfmathsetmacro{\ph}{79.1066}
\pgfmathsetmacro{\ps}{26.3689}
\pgfmathsetmacro{\r}{0.881917}
\pgfmathsetmacro{\h}{0.790156}
\pgfmathsetmacro{\al}{0.623490}
\draw[gray] (0,0) circle (1);
\draw[->,gray] (-1.15,0) -- (1.25,0) node[right] {$x$};
\foreach \k in {0,...,6}{\fill[cnew!70] ({cos(360*\k/7)},{sin(360*\k/7)}) circle (0.5pt);}
\draw[cnew!50,dashed] (1,0) \foreach \k in {1,...,6}{-- ({cos(360*\k/7)},{sin(360*\k/7)})} -- cycle;
\draw[cold,thick] (0,0) -- ({1.12*cos(\ph)},{1.12*sin(\ph)});
\draw[cold,thin] ({cos(\ph)},{sin(\ph)}) -- ({cos(\ph)},0);
\draw[cold] (0.16,0) arc (0:\ph:0.16);
\node[cold] at (0.09,0.22) {$\varphi$};
\draw[cnew,thick] (0,0) -- ({1.12*cos(\ps)},{1.12*sin(\ps)});
\draw[cnew] (0.3,0) arc (0:\ps:0.3);
\node[cnew] at (0.36,0.08) {$\psi$};
\fill[cnew] ({\r*cos(\ps)},{\r*sin(\ps)}) circle (0.6pt) node[above left=-1pt] {$M$};
\draw[cnew,thick] ({\r*cos(\ps)},{\r*sin(\ps)}) -- (\h,0);
\fill[cnew] (\h,0) circle (0.6pt) node[below right=-1pt] {$H$};
\fill[cnew] (\al,0) circle (0.6pt) node[below=1pt] {$\alpha$};
\draw[cnew,thick] (\al,0) -- (\al,{sin(360/7)+0.08});
\fill[cnew] ({cos(360/7)},{sin(360/7)}) circle (0.8pt) node[above right=-1pt] {$P_1$};
\fill[cold] (1,0) circle (0.6pt) node[below right=-1pt] {$P_0$};
\fill (0,0) circle (0.6pt) node[below left=-1pt] {$O$};
\end{tikzpicture}
\figcap{Слева: угол $\varphi$ строится по своему косинусу: в прямоугольном треугольнике с катетами $AB=1$ и
$AC=3\sqrt3$ гипотенуза $BC=2\sqrt7$, так что $\cos\varphi=\frac{AB}{BC}=\frac1{2\sqrt7}$. Справа: построение
правильного семиугольника трисектором. Угол $\varphi$ (синий прямоугольный треугольник справа подобен левому) делится на три;
$OM=\frac{\sqrt7}3$ на стороне угла $\psi=\frac\varphi3$, проекция $OH=\frac{\sqrt7}3\cos\psi$, и $\alpha=OH-\frac16=\cos\frac{2\pi}7$.}
\end{keepfig}
\end{center}

\emph{Замечания.} Дискриминант $\mu$ не меняется при сдвиге и равен $(\frac78)^2$~--- квадрату в $\Q$; поэтому группа Галуа
$\mu$ есть $A_3\cong\Z/3\Z$, и все корни $\mu$ лежат в $\Q[\alpha]$: $\cos\frac{4\pi}7=2\alpha^2-1$,
$\cos\frac{6\pi}7=4\alpha^3-3\alpha$. Правильный $11$-угольник трисектором не построить (\rf{1.42a}, \rf{5.13a}). Вообще, по
теореме Глисона (приводим без доказательства) правильный $n$-угольник строится циркулем, линейкой и трисектором тогда
и только тогда, когда $n=2^r3^sp_1\cdots p_k$, где $p_i>3$~--- различные простые вида $2^u3^v+1$; например, $7=2\cdot3+1$ и
$13=2^2\cdot3+1$, а $11$ не подходит.

\begin{exercise}[1-5]
Пусть $f(X)$~--- неприводимый многочлен степени $n$ над $F$, а $E$~--- расширение $F$ степени $[E:F]=m$. Покажите, что
если $\gcd(m,n)=1$, то $f$ неприводим над $E$.
\end{exercise}

\noindent\emph{Решение.} Пусть $g$~--- унитарный неприводимый множитель $f$ в $E[X]$, и $L=E[X]/(g)=E[\alpha]$~--- его поле
корня (\rf{1.25}); здесь $g(\alpha)=0$, $[L:E]=\deg g$. Так как $f(\alpha)=0$ и $f$ неприводим над $F$, минимальный многочлен
$\alpha$ над $F$~--- это $f$ (с точностью до множителя), и $[F[\alpha]:F]=n$. Рассмотрим две башни внутри $L$:
\[
F\subset E\subset L,\qquad F\subset F[\alpha]\subset L .
\]
По \rf{1.20} $[L:F]=m\deg g$, и $n=[F[\alpha]:F]$ делит $[L:F]$. Значит, $n\mid m\deg g$, а так как $\gcd(m,n)=1$, то
$n\mid\deg g$. Но $g\mid f$, т.\,е.\ $\deg g\le n$. Поэтому $\deg g=n$, $g$ совпадает с $f$ с точностью до множителя, и $f$
неприводим над $E$.

\begin{center}
\begin{keepfig}
\centering
\begin{tikzpicture}[every node/.style={font=\small}]
\node (F) at (0,0) {$F$};
\node (E) at (-1.3,1.2) {$E$};
\node (Fa) at (1.3,1.2) {$F[\alpha]$};
\node (L) at (0,2.4) {$L=E[\alpha]$};
\draw (F) -- node[below left=-2pt] {$m$} (E);
\draw (F) -- node[below right=-2pt] {$n$} (Fa);
\draw (E) -- node[above left=-2pt,cnew] {$\deg g$} (L);
\draw (Fa) -- (L);
\end{tikzpicture}
\figcap{$n$ делит $[L:F]=m\deg g$, и из $\gcd(m,n)=1$ следует $n\mid\deg g$.}
\end{keepfig}
\end{center}

\emph{Замечание.} Утверждение означает, что кольцо $E[X]/(f)$~--- поле (\rf{1.25} и абзац после построения
$F[X]/(f)$), или, что то же самое по \rf{1.23}, что в нём нет делителей нуля. Доказывать это непосредственно~---
например, раскладывая коэффициенты возможных множителей $f$ в $E[X]$ по базису $E$ над $F$,~--- неудобно: условия на
коэффициенты получаются громоздкими. Удобнее, как выше, перейти к полю корня неприводимого множителя $g$ и сравнить
степени. То же рассуждение без предположения $\gcd(m,n)=1$ даёт: степень каждого неприводимого множителя $f$ в $E[X]$
делится на $\frac n{\gcd(m,n)}$ (из $n\mid m\deg g$). Например, $X^4-2$ над $\Q(\sqrt2)$ раскладывается в произведение
$(X^2-\sqrt2)(X^2+\sqrt2)$ неприводимых множителей степени $2=\frac42$.

\begin{exercise}[1-6]
Покажите, что не существует многочлена $f(X)\in\Z[X]$ степени $>1$, неприводимого по модулю $p$ для всех простых $p$.
\end{exercise}

\noindent\emph{Решение.} Для унитарного $f$ степени $n\ge2$ рассуждение совсем короткое: значения $\pm1$ многочлен $f$
принимает лишь в конечном числе точек, поэтому $f(m)\ne\pm1$ для некоторого $m\in\Z$; возьмём простое $p\mid f(m)$ (при
$f(m)=0$~--- любое). Тогда $\bar m$~--- корень приведения $\bar f\in\F_p[X]$, которое имеет ту же степень $n\ge2$ (старший
коэффициент равен $1$), так что $\bar f$ приводимо. В общем случае простое $p\mid f(m)$ может делить старший коэффициент,
и рассуждаем так.

Пусть $f=a_nX^n+\dots+a_0$, $n\ge2$. Достаточно найти простое $p\nmid a_n$, по модулю которого $f$
имеет корень: тогда приведение $\bar f\in\F_p[X]$ имеет ту же степень $n\ge2$ и делится на $X-\bar r$ (1-3(b)), т.\,е.\
приводимо. Покажем, что простых чисел $p$, по модулю которых $f$ имеет корень (т.\,е.\ $p\mid f(r)$ для некоторого
$r\in\Z$), бесконечно много; тогда среди них есть и не делящие $a_n$.

Если $a_0=0$, то $f(p)\equiv0\pmod p$ при всех $p$. Пусть $a_0\ne0$, и пусть таких простых конечное число: $p_1,\dots,p_s$.
Положим $Q=p_1\cdots p_s$. Тогда
\[
f(a_0Qx)=a_0\bigl(1+Qx\,h(x)\bigr),\qquad h\in\Z[x],
\]
так как каждый член $a_k(a_0Qx)^k$ с $k\ge1$ делится на $a_0Qx$. Многочлен $1+Qxh(x)$ имеет степень $n\ge1$, поэтому
значения $\pm1$ он принимает лишь в конечном числе точек, и найдётся целое $x$, при котором $N=1+Qxh(x)\ne\pm1$. Пусть
$q$~--- простой делитель $N$. Тогда $q\mid f(a_0Qx)$, т.\,е.\ $f$ имеет корень по модулю $q$, но $N\equiv1\pmod{p_i}$ для всех
$i$, так что $q\ne p_i$~--- противоречие.

Обратное явление тоже бывает: многочлен $X^4-10X^2+1$ неприводим над $\Q$, но приводим по модулю каждого простого (\rf{1.18c};
ср.\ 1-2), а многочлен из упражнения 1-8 имеет корень по модулю каждого простого, хотя не имеет рациональных корней.

\begin{exercise}[1-7]
Пусть $\alpha=\sqrt[3]2$, и $R$~--- множество комплексных чисел вида $a+b\alpha+c\alpha^2$, $a,b,c\in\Q$. Покажите, что $R$~---
поле.
\end{exercise}

\noindent\emph{Решение.} Многочлен $X^3-2$ неприводим над $\Q$ (Эйзенштейн, $p=2$), поэтому это минимальный многочлен
$\alpha$, и $\Q[\alpha]\cong\Q[X]/(X^3-2)$~--- поле (\rf{1.25}), каждый элемент которого однозначно записывается в виде
$a+b\alpha+c\alpha^2$. Значит, $R=\Q[\alpha]$~--- поле. (Что $R$~--- кольцо, видно и непосредственно: при умножении
степени $\alpha^3=2$, $\alpha^4=2\alpha$ понижаются.) Второй способ: $R$~--- подкольцо $\CC$, т.\,е.\ область
целостности, конечномерная над $\Q$, и по \rf{1.23} это поле.

Обратный элемент можно указать явно. Прямое перемножение (с $\alpha^3=2$) проверяет тождество
\[
(a+b\alpha+c\alpha^2)\bigl((a^2-2bc)+(2c^2-ab)\alpha+(b^2-ac)\alpha^2\bigr)=a^3+2b^3+4c^3-6abc .
\]
Правая часть $N(a,b,c)$ не равна нулю при $(a,b,c)\ne(0,0,0)$: она равна произведению
$\prod_{j=0}^2(a+b\omega^j\alpha+c\omega^{2j}\alpha^2)$, $\omega=e^{2\pi i/3}$, по всем трём корням $\omega^j\alpha$ многочлена $X^3-2$, а
каждый множитель ненулевой, так как корень неприводимого многочлена степени $3$ не может быть корнем ненулевого
многочлена $a+bX+cX^2$ степени $\le2$. Поэтому
\[
(a+b\alpha+c\alpha^2)^{-1}=\frac{(a^2-2bc)+(2c^2-ab)\alpha+(b^2-ac)\alpha^2}{a^3+2b^3+4c^3-6abc}.
\]
Например, $(1+\alpha)^{-1}=\frac{1-\alpha+\alpha^2}3$ (проверка: $(1+\alpha)(1-\alpha+\alpha^2)=1+\alpha^3=3$).

Найти обратный элемент можно и не угадывая тождество, двумя систематическими способами.

(i) \emph{Алгоритм Евклида.} Пусть $g=cX^2+bX+a\ne0$. Многочлены $X^3-2$ и $g$ взаимно просты в $\Q[X]$: их общий
делитель делит неприводимый $X^3-2$, а сам $X^3-2$ не делит $g$ (степень меньше). По \rf{1.8} существуют $u,v\in\Q[X]$ с
$u(X)(X^3-2)+v(X)g(X)=1$, и подстановка $X=\alpha$ даёт $v(\alpha)g(\alpha)=1$, т.\,е.\ $(a+b\alpha+c\alpha^2)^{-1}=v(\alpha)$.
Например, для $g=X+1$ один шаг деления даёт $X^3-2=(X+1)(X^2-X+1)-3$, откуда $(X+1)\cdot\frac{X^2-X+1}3-\frac{X^3-2}3=1$ и
снова $(1+\alpha)^{-1}=\frac{1-\alpha+\alpha^2}3$.

(ii) \emph{Линейная система.} Ищем $x+y\alpha+z\alpha^2$ с $(a+b\alpha+c\alpha^2)(x+y\alpha+z\alpha^2)=1$. Раскрыв скобки
и заменив $\alpha^3=2$, $\alpha^4=2\alpha$, получаем систему
\[
\begin{pmatrix}a&2c&2b\\ b&a&2c\\ c&b&a\end{pmatrix}\begin{pmatrix}x\\ y\\ z\end{pmatrix}=\begin{pmatrix}1\\0\\0\end{pmatrix};
\]
её матрица~--- матрица умножения на $\beta=a+b\alpha+c\alpha^2$ в базисе $1,\alpha,\alpha^2$, а определитель равен
$a^3+2b^3+4c^3-6abc=N(a,b,c)$ (это норма $\beta$, ср.\ \rf{5.42b}). Он не равен нулю при $(a,b,c)\ne(0,0,0)$: иначе
нашёлся бы ненулевой столбец $(x,y,z)$ с нулевым произведением, т.\,е.\ $\beta\gamma=0$ для $\gamma=x+y\alpha+z\alpha^2\ne0$
(элементы $1,\alpha,\alpha^2$ линейно независимы), откуда $\beta=0$ в $\CC$; тогда $\alpha$ был бы корнем ненулевого
многочлена $a+bX+cX^2$ и имел бы степень $<3$. По правилу Крамера решение системы~--- в точности формула для
$(a+b\alpha+c\alpha^2)^{-1}$, выписанная выше.

\begin{exercise}[1-8]
С помощью свойств символа Лежандра покажите, что многочлен $(X^2-13)(X^2-17)(X^2-221)$ имеет корни по модулю
каждого целого числа (но не в $\Z$).
\end{exercise}

\noindent\emph{Решение.} Пусть $f=(X^2-13)(X^2-17)(X^2-221)$; заметим, что $221=13\cdot17$. Рациональных корней у $f$ нет:
$13$, $17$, $221$ не являются квадратами целых чисел, а значит, и рациональных (\rf{1.11}). Требуется доказать, что для
каждого $m\ge1$ сравнение $f(x)\equiv0\pmod m$ разрешимо. Напомним, что для нечётного простого $p$ и $a\not\equiv0\pmod p$
символ Лежандра $\bigl(\frac ap\bigr)$ равен $1$, если $a$~--- квадрат по модулю $p$, и $-1$ иначе, и он мультипликативен:
$\bigl(\frac{ab}p\bigr)=\bigl(\frac ap\bigr)\bigl(\frac bp\bigr)$ (это \rf{1.18b}: квадраты~--- подгруппа индекса $2$ в $\F_p^\times$).
Равносильное определение~--- \emph{критерий Эйлера}: $\bigl(\frac ap\bigr)\equiv a^{(p-1)/2}\pmod p$. Действительно, если
$\F_p^\times=\langle g\rangle$ (упражнение 1-3(e)) и $a=g^k$, то $a$~--- квадрат тогда и только тогда, когда $k$ чётно
($p-1$ чётно), а $a^{(p-1)/2}=\bigl(g^{(p-1)/2}\bigr)^k=(-1)^k$, так как $g^{(p-1)/2}$~--- элемент порядка $2$, т.\,е.\ корень
$X^2-1$, отличный от $1$.

\emph{Простые модули.} При $p=2$: $X^2-13\equiv X^2-1$, корень $1$. При $p=13$: $17\equiv4=2^2$, корень $2$ у
$X^2-17$. При $p=17$: $13\equiv64=8^2$, корень $8$ у $X^2-13$. Пусть $p\ne2,13,17$. Если $\bigl(\frac{13}p\bigr)=\bigl(\frac{17}p\bigr)=-1$,
то $\bigl(\frac{221}p\bigr)=(-1)(-1)=1$. Значит, хотя бы одно из чисел $13,17,221$~--- ненулевой квадрат по модулю $p$, и
соответствующий множитель имеет корень. (Квадратичный закон взаимности \rf{5.11l} здесь не нужен; с ним можно
сказать точнее, какой множитель работает: $\bigl(\frac{13}p\bigr)=\bigl(\frac p{13}\bigr)$ и
$\bigl(\frac{17}p\bigr)=\bigl(\frac p{17}\bigr)$, так как показатели $\frac{13-1}2\cdot\frac{p-1}2$ и $\frac{17-1}2\cdot\frac{p-1}2$
чётны.)

\emph{Степени простых.} Для нечётного $p$ корень по модулю $p$ поднимается до корня по модулю любой степени $p$ по
лемме Гензеля \rf{1.18g}: для $g=X^2-a$ её условие $p\nmid g'(x_1)=2x_1$ выполнено, если $p\nmid a$. Поэтому для $p=13$ и
$p=17$ берём $a=17$ и $a=13$ соответственно (очевидный корень $0$ множителя $X^2-13$ по модулю $13$ не поднимается,
см.\ \rf{1.18g}), а для прочих нечётных $p$~--- то из чисел $13,17,221$, которое является квадратом по модулю $p$.

Для $p=2$ нужен множитель $X^2-17$: число $17\equiv1\pmod8$ является квадратом по модулю каждого $2^k$. Действительно,
$1^2\equiv17\pmod{2^k}$ при $k\le3$, и если $x^2\equiv17\pmod{2^k}$, $k\ge3$, $x$ нечётно, то либо $x^2\equiv17\pmod{2^{k+1}}$,
либо $x^2\equiv17+2^k\pmod{2^{k+1}}$, и тогда
\[
(x+2^{k-1})^2=x^2+2^kx+2^{2k-2}\equiv x^2+2^k\equiv17\pmod{2^{k+1}},
\]
так как $x$ нечётно и $2k-2\ge k+1$. (Числа $13\equiv221\equiv5\pmod8$ не являются квадратами уже по модулю $8$: квадраты
нечётных чисел сравнимы с $1$ по модулю $8$. Поэтому множитель $X^2-17$ здесь действительно нужен.) Лемма \rf{1.18g}
здесь неприменима ($g'(x)=2x$ чётно), а сделанный шаг~--- частный случай её <<сильной>> формы, которую мы не доказываем:
если $g(x_1)\equiv0\pmod{p^{2e+1}}$, где $p^e$~--- точная степень $p$, делящая $g'(x_1)$, то корни $g$ по модулю всех
степеней $p$ существуют (здесь $g=X^2-17$, $x_1=1$, $g(1)=-16$, $e=1$).

\emph{Произвольный модуль.} Пусть $m=\prod p^{k_p}$. Для каждого $p\mid m$ выберем, как выше, множитель $X^2-a_p$ многочлена
$f$ и $r_p$ с $r_p^2\equiv a_p\pmod{p^{k_p}}$. По китайской теореме об остатках найдётся $x$ с $x\equiv r_p\pmod{p^{k_p}}$ для
всех $p\mid m$. Тогда $f(x)\equiv0\pmod{p^{k_p}}$ для каждого $p$ (один из множителей $f(x)$ делится на $p^{k_p}$), и,
следовательно, $f(x)\equiv0\pmod m$. Например, при $m=8\cdot13\cdot17=1768$ берём
$x\equiv1\pmod8$ (корень $X^2-17$), $x\equiv2\pmod{13}$ (корень $X^2-17$), $x\equiv8\pmod{17}$ (корень $X^2-13$) и
получаем $x=977$: число $977^2-17$ делится на $8\cdot13$, а $977^2-13$~--- на $17$, так что $f(977)\equiv0\pmod{1768}$.

Ни один множитель в отдельности не годится для всех $m$: $x^2\equiv13$ неразрешимо по модулю $5$, а $x^2\equiv17$ и
$x^2\equiv221$~--- по модулю $3$; более того, по модулю $13^k$ ($k\ge2$) годится только $X^2-17$, а по модулю $17^k$
($k\ge2$)~--- только $X^2-13$ (для $2^k$ см.\ выше). Существенны и знаки: для $h=(X^2+13)(X^2+17)(X^2+221)$ утверждение
неверно. По модулю $43$: $13\equiv20^2$, $17\equiv19^2$, а $-1$~--- не квадрат (по критерию Эйлера $(-1)^{21}=-1$),
поэтому $-13$, $-17$, $-221$~--- не квадраты, и $h$ не имеет корней по модулю $43$. Нет у $h$ корней и по модулю $16$:
при чётном $x$ все три множителя нечётны, а при нечётном $x$ имеем $x^2\equiv1\pmod8$, множители сравнимы с $6,2,6$
по модулю $8$, и их произведение делится на $8$, но не на $16$.

Итак, $f$ имеет корень по модулю каждого $m$, но не имеет корней в $\Z$ (и в $\Q$). Это показывает, что для многочленов
нет <<локально-глобального принципа>>: разрешимость по всем модулям не влечёт разрешимости в целых числах. Для одного
многочлена $X^2-c$ так не бывает: целое $c$, являющееся квадратом по модулю каждого $m$, есть квадрат целого числа
(этот более глубокий факт мы не доказываем).
